\begin{abstract}
Prompts such as ``most underrated X'' are widely used for discovery, but it is unclear whether they reveal genuine model novelty or just recycled crowd-approved picks. We introduce a reproducible diagnostic that combines real API outputs, a data-grounded underratedness proxy from \movielens, and social-consensus overlap signals from \reddit. We evaluate two model families (\gptfourone and \gptfive), nine movie genres, three prompt conditions (baseline, anti-consensus, anti-consensus+reasoned), and two repeated runs per cell, yielding 108 model outputs. Our robustness metric, the \uns hit score, combines percentile underratedness with a direct title-hit overlap penalty. For \gptfourone, anti-consensus prompting substantially reduces overlap-hit rate from 38.9\% in baseline to 22.2\% (anti-consensus) and 11.1\% (anti-consensus+reasoned), with the best mean \uns hit score under anti-consensus+reasoned (0.669 vs.\ 0.562 baseline). However, paired permutation tests with Holm correction show no statistically significant gains after multiple-comparison adjustment. For \gptfive in this setup, anti-consensus variants suffer severe structured-output failures, which dominate metric outcomes. These results suggest that ``most underrated'' prompts are promising as novelty probes, but conclusions depend strongly on metric construction quality and output-validation reliability. We release a modular pipeline and report limitations to support stronger future benchmark design.
\end{abstract}
